\section{Gemma con aprendizaje en contexto}
\label{chap:vlm}

\subsection{Modelo y servicio}

Se utiliza \texttt{google/gemma-4-E4B-it} \cite{gemma4card}, una variante preparada para seguir instrucciones que recibe conjuntamente las imágenes de ECG y el texto de la tarea. Sus pesos permanecen fijos durante los experimentos, de modo que cada predicción se obtiene a partir de las instrucciones y de los ejemplos incluidos en la consulta. El número y el origen de esas demostraciones permiten estudiar cómo cambia su respuesta al disponer de distintos contextos.

El modelo se ejecuta en un clúster mediante vLLM, que mantiene el servicio de inferencia encargado de generar las respuestas. Dentro de cada tarea se mantienen el mismo modelo y los parámetros de inferencia, incluida una temperatura configurada de cero para reducir la variación asociada a la generación de texto. La Tabla~\ref{tab:config_vlm_appendix} recoge la configuración completa y el capítulo~\ref{chap:gestion_etica} desarrolla las implicaciones de privacidad asociadas al uso de grandes modelos de lenguaje con datos clínicos.

\subsection{Prompt y respuesta}

El prompt es el mensaje de entrada que recibe el modelo. Tras seleccionar las demostraciones según la sección~\ref{sec:condiciones_vlm}, se construye con la instrucción de sistema, los ejemplos y la imagen de consulta, en ese orden. Cada ejemplo incluye una imagen V1 y su respuesta en JSON, un formato que organiza la información en campos y valores. La consulta exige el mismo esquema. La Figura~\ref{fig:icl_context_window} representa la secuencia para el caso de clasificación de datos simulados.

Cada demostración se envía como una entrada de usuario con la imagen y una respuesta de asistente con la etiqueta conocida. En este estudio, no se han seleccionado demostraciones por similitud visual con la consulta ni se han estudiado permutaciones del orden, conteniendo así la última entrada la imagen que debe ser clasificada.

\begin{figure}[H]
  \centering
  % Fuente editable: assets/results/vlm_icl_context_window.svg.
  \safeincludegraphics[width=\textwidth]{results/vlm_icl_context_window.pdf}
  \caption{Consulta con demostraciones y predicción de los tres rasgos del ECG simulado mediante ICL.}
  \label{fig:icl_context_window}
\end{figure}

La tarea del simulador solicita los tres rasgos representados en la figura y la del conjunto cedido por el HUCA, una etiqueta positiva o negativa.

El formato de salida vincula la respuesta del modelo con las etiquetas utilizadas en la evaluación. En el simulador se conserva la decisión sobre cada rasgo y se obtiene una clasificación positiva derivada cuando los tres están presentes. En los registros provenientes del HUCA, la respuesta expresa directamente la etiqueta de Brugada solicitada en la tarea. El apéndice~\ref{app:prompts} recoge las instrucciones y los formatos empleados en cada experimento.
